\chapter{Cystic Fibrosis}
\index{Cystic Fibrosis}

\section*{Definition and Overview}
Cystic fibrosis (CF) is an inherited condition that affects the cells that produce mucus, sweat, and digestive juices. A faulty gene causes the body to produce thick, sticky mucus that clogs the lungs and blocks ducts in the pancreas and other organs. This leads to repeated lung infections, progressive lung damage, and problems digesting food. Cystic fibrosis is the most common life-limiting inherited condition in people of European descent, but treatments have improved dramatically, and many people now live into their forties, fifties, or longer. Modern care involves daily treatments, physiotherapy, enzyme supplements, and, in recent years, highly effective modulator medicines that target the underlying defect.

\section*{Epidemiology}
Cystic fibrosis affects about 1 in 2,500--3,500 newborns of European ancestry, with around one in 25 people carrying the faulty gene. In many countries, about one baby in 2,500 is born with CF, and it is detected soon after birth through newborn screening. The condition affects boys and girls equally. Survival has improved steadily: the median predicted survival in many countries is now around 50 years or more, a dramatic change from the past, driven by specialist care, new medicines, and lung transplantation.

\section*{Aetiology and Pathophysiology}
Cystic fibrosis is caused by mutations in the CFTR gene, which controls the movement of salt and water in and out of cells.

\begin{itemize}
    \item \textbf{CFTR gene:} mutations in the cystic fibrosis transmembrane conductance regulator gene
    \item \textbf{Inheritance:} autosomal recessive, requiring two faulty copies, one from each parent
    \item \textbf{Most common mutation:} F508del, found in most people with CF
    \item \textbf{Salt and water defect:} abnormal chloride transport makes secretions thick and sticky
    \item \textbf{Lung disease:} thick mucus traps bacteria, causing chronic infection, inflammation, and scarring
    \item \textbf{Pancreatic disease:} blocked ducts prevent digestive enzymes reaching the gut, causing malabsorption
    \item \textbf{Other organs:} the liver, sinuses, sweat glands, and reproductive organs are affected
    \item \textbf{Modulator medicines:} newer drugs correct or enhance the function of the faulty protein
\end{itemize}

\section*{Clinical Features}
Symptoms vary with age and severity.

\begin{itemize}
    \item \textbf{Lungs:} persistent cough, frequent chest infections, wheeze, and shortness of breath
    \item \textbf{Digestion:} poor weight gain, greasy stools, and vitamin deficiencies
    \item \textbf{Sinuses:} chronic sinusitis and nasal polyps
    \item \textbf{Salt loss:} salty-tasting skin and salt depletion in heat
    \item \textbf{Diabetes:} CF-related diabetes in later life
    \item \textbf{Liver:} liver disease in some people
    \item \textbf{Fertility:} most men with CF are infertile; women have reduced fertility
    \item \textbf{Clubbing:} widening of the fingertips from chronic lung disease
\end{itemize}

\section*{Differential Diagnosis}
Before diagnosis, other causes of the features are considered, though newborn screening now identifies most cases early.

\begin{itemize}
    \item Asthma and other causes of chronic cough and wheeze
    \item Recurrent pneumonia or bronchiectasis from other causes
    \item Coeliac disease and other causes of malabsorption and poor growth
    \item Primary ciliary dyskinesia, a different inherited condition
    \item Immunodeficiency causing recurrent infections
    \item Chronic pancreatitis from other causes
\end{itemize}

\section*{When to Seek Medical Care}
People with cystic fibrosis are cared for by specialist CF clinics and follow a treatment plan. Contact the CF team promptly for any new or worsening symptoms, especially fever, increased cough, or shortness of breath.

\textbf{Contact your local emergency services for:}
\begin{itemize}
    \item Severe difficulty breathing or a sudden worsening of lung symptoms
    \item Coughing up significant amounts of blood
    \item Severe chest or abdominal pain
    \item Signs of a serious infection: high fever, confusion, or collapse
    \item An inability to eat or drink with persistent vomiting
\end{itemize}

\section*{Diagnosis and Investigations}
Diagnosis is confirmed by sweat testing and genetic testing, with monitoring throughout life.

\begin{itemize}
    \item \textbf{Newborn screening:} a heel-prick blood test, followed by confirmation
    \item \textbf{Sweat test:} measures chloride in sweat; high levels confirm CF
    \item \textbf{Genetic testing:} identifies the specific CFTR mutations
    \item \textbf{Chest X-ray and CT:} assessment of lung damage and infection
    \item \textbf{Lung function tests:} spirometry to monitor breathing capacity
    \item \textbf{Sputum cultures:} to identify lung infections
    \item \textbf{Pancreatic function and nutrition:} stool tests, weight, and vitamin levels
    \item \textbf{Glucose tolerance testing:} for CF-related diabetes
\end{itemize}

\section*{Management}
CF care is delivered by a multidisciplinary team and includes daily treatment.

\begin{itemize}
    \item \textbf{Chest physiotherapy:} daily airway clearance techniques and exercise
    \item \textbf{Inhaled medicines:} mucus-thinning agents, bronchodilators, and antibiotics
    \item \textbf{Antibiotics:} oral, inhaled, or intravenous treatment of infections
    \item \textbf{Pancreatic enzymes:} taken with food to digest nutrients
    \item \textbf{Nutritional support:} high-energy diet, vitamin supplements, and feeding support
    \item \textbf{CFTR modulator medicines:} highly effective drugs that treat the underlying defect, dramatically improving lung function and quality of life
    \item \textbf{CF-related diabetes:} insulin and diet management
    \item \textbf{Liver care:} monitoring and treatment of liver disease
    \item \textbf{Lung transplantation:} for advanced disease
    \item \textbf{Psychological and social support:} for the person and family
\end{itemize}

\section*{Complications}
\begin{itemize}
    \item Progressive lung damage and bronchiectasis
    \item Chronic infection with resistant bacteria, including \textit{Pseudomonas}
    \item Respiratory failure
    \item Malnutrition and vitamin deficiencies
    \item CF-related diabetes
    \item Liver disease and gallstones
    \item Bone disease and osteoporosis
    \item Male infertility and reduced female fertility
    \item Complications of transplantation
    \item Mental health problems, including anxiety and depression
\end{itemize}

\section*{Prognosis and Outcomes}
The outlook for cystic fibrosis has transformed over recent decades. Median survival in many countries is now around 50 years, and children born today can expect to live longer still. CFTR modulator therapies have produced remarkable improvements in lung function, weight, and quality of life for eligible people, and research continues. The disease remains serious and progressive, requiring daily treatment and specialist care, but most adults with CF complete education, work, and have fulfilling lives.

\section*{Prevention and Self-Care}
\begin{itemize}
    \item Attend a specialist CF clinic and follow the treatment plan every day
    \item Perform airway clearance and take all inhaled and oral medicines
    \item Take enzymes with every meal and snack, and follow dietary advice
    \item Exercise regularly to maintain lung health
    \item Stay up to date with vaccinations, including influenza, pneumococcal, and COVID-19
    \item Avoid smoke and respiratory infections; practise good hand hygiene
    \item Keep good blood sugar control if you have CF-related diabetes
    \item Discuss genetic testing and family planning options with the team
    \item Access mental health and peer support
\end{itemize}

\section*{Sources and Further Reading}
The following reputable sources provide further information for healthcare students and patients seeking reliable, current advice about cystic fibrosis.

\begin{itemize}
    \item Cystic Fibrosis Australia. \textit{Information and support.} https://www.cysticfibrosis.org.au/
    \item Cystic Fibrosis Foundation. \textit{Cystic fibrosis information.} https://www.cff.org/
    \item Healthdirect Australia. \textit{Cystic fibrosis.} https://www.healthdirect.gov.au/
    \item NHS. \textit{Cystic fibrosis.} https://www.nhs.uk/conditions/cystic-fibrosis/
    \item Lung Foundation Australia. \textit{Cystic fibrosis resources.} https://lungfoundation.com.au/
    \item MedlinePlus. \textit{Cystic fibrosis.} https://medlineplus.gov/ency/article/000107.htm
\end{itemize}
